\documentclass[a4paper,twoside,10pt]{article}

% --- LAYOUT ---
\usepackage{geometry}
\geometry{a4paper, left=22mm, right=22mm, top=18mm, bottom=18mm}
\usepackage{parskip}
\setlength{\parskip}{4pt}
\usepackage{setspace}
\setstretch{1.1}

% --- WISKUNDE ---
\usepackage{amsmath, amssymb, amsfonts, amsthm}
\usepackage{mathrsfs}
\usepackage{bm}
\usepackage{esint}
\usepackage{siunitx}

% --- TABELLEN ---
\usepackage{booktabs}
\usepackage{array}

% --- OVERIGE ---
\usepackage{graphicx}
\usepackage[hidelinks]{hyperref}
\usepackage[T1]{fontenc}
\usepackage[utf8]{inputenc}
\usepackage{multicol}

% --- THEOREM STIJLEN ---
\theoremstyle{definition}
\newtheorem{postulate}{Postulaat}[section]
\newtheorem{derivation}{Afleiding}[section]

% --- COMMANDO'S ---
\newcommand{\dbar}{{\mkern+3mu\mathchar'26\mkern-12mu d}}
% \pd[const]{teller}{noemer}  =>  (\partial teller / \partial noemer)_const
\newcommand{\pd}[3][]{\left(\dfrac{\partial #2}{\partial #3}\right)_{\!#1}}

% --- TITEL ---
\title{\Large\bfseries Samenvatting Thermische Fysica}
\author{T. Cools}
\date{Prof.\ dr.\ N. Jachowicz (2025--2026)}

\begin{document}
\maketitle
\tableofcontents
\bigskip\hrule\bigskip

%%=============================================================
\section{Grondbeginselen (H1--H3)}
%%=============================================================

\subsection{Systemen, Processen en Wanden}

\textbf{Extensief}: evenredig met massa ($V, U, S, \ldots$).\quad
\textbf{Intensief}: onafhankelijk van massa ($P, T, \ldots$).

\begin{center}
\renewcommand{\arraystretch}{1.2}
\begin{tabular}{@{}ll@{}}
\toprule
Concept & Definitie \\
\midrule
Geïsoleerd systeem  & geen energie- of stofuitwisseling met omgeving \\
Gesloten systeem    & energie-uitwisseling, geen stofuitwisseling \\
Diathermische wand  & doorlaatbaar voor warmte \\
Adiabatische wand   & niet doorlaatbaar voor warmte \\
Therm.\ evenwicht   & mechanisch + thermisch + chemisch evenwicht simultaan \\
Quasistatisch proc. & verloopt zo traag dat systeem steeds in evenwicht is \\
\bottomrule
\end{tabular}
\end{center}

\subsection{Nulde Wet en Temperatuurschalen}

\begin{postulate}[Nulde wet]
Als A in thermisch evenwicht is met C, en B ook met C,
dan zijn A en B in thermisch evenwicht met elkaar.
\end{postulate}

Kelvinschaal via tripelpunt van water
($T_{tp}=273{,}16\,\mathrm{K}$, $P_{tp}=611{,}73\,\mathrm{Pa}$):
\[
  T = (273{,}16\,\mathrm{K})\,\frac{P}{P_{tp}},
  \qquad
  T(\mathrm{K}) = T({}^{\circ}\mathrm{C}) + 273{,}15
\]

\subsection{Toestandsvergelijkingen en Sleutelgrootheden}

Toestandsvergelijking: $f(P,V,T)=0$ — slechts 2 van de 3 variabelen zijn onafhankelijk.

\[
  \beta = \frac{1}{V}\pd[P]{V}{T}
  \quad\text{(vol.-uitzettingscoëff.)},
  \qquad
  \kappa = -\frac{1}{V}\pd[T]{V}{P}
  \quad\text{(isoth.\ drukbaarheid)}
\]
\[
  dV = V\beta\,dT - V\kappa\,dP, \qquad dP = \frac{\beta}{\kappa}\,dT - \frac{1}{V\kappa}\,dV
\]

\textbf{Ideaal gas}: $PV = nRT$.\quad
Viriaalontwikkeling: $Z = PV/nRT = 1 + B(T)/V + C(T)/V^2 + \cdots$

\textbf{Overzicht veralgemeende systemen:}
\begin{center}
\renewcommand{\arraystretch}{1.2}
\begin{tabular}{@{}lllll@{}}
\toprule
Systeem & Ext.\ coord.\ $X$ & Int.\ kracht $Y$ & $\dbar W$ & Toestandsvgl. \\
\midrule
Hydrostatisch   & $V$   & $-P$      & $-P\,dV$        & $PV = nRT$ \\
Gespannen snaar & $L$   & $F$       & $F\,dL$         & $F=k(L-L_0)$;\; $\beta=3\alpha$ \\
Oppervlaktefilm & $A$   & $\gamma$  & $\gamma\,dA$    & $(\gamma-\gamma_w)A = cT$ \\
Diëlektricum    & $\Pi$ & $E$       & $E\,d\Pi$       & $\Pi/V = a + b/T$ \\
Paramagneticum  & $M$   & $\mu_0 H$ & $-\mu_0 H\,dM$ & $M = (C_c/T)\,H$ (Curie) \\
\bottomrule
\end{tabular}
\end{center}

%%=============================================================
\section{Arbeid, Warmte en Eerste Wet (H4--H5)}
%%=============================================================

\subsection{Arbeid}

Arbeid is \textbf{padafhankelijk} ($\dbar W \neq dW_{\text{exact}}$).\quad
Algemeen: $\dbar W = \sum_i Y_i\,dX_i$.

Quasistatisch hydrostatisch: $W = -\displaystyle\int_{V_i}^{V_f} P\,dV$.

\begin{center}
\renewcommand{\arraystretch}{1.2}
\begin{tabular}{@{}lll@{}}
\toprule
Proces & Formule & Opmerking \\
\midrule
Isobaar              & $W = -P(V_f-V_i)$                        & \\
Isochoor             & $W = 0$                                   & \\
Isotherm, ideaal gas & $W = -nRT\ln(V_f/V_i)$                   & \\
Isotherm, vaste stof & $W \approx \frac{\kappa V}{2}(P_f^2-P_i^2)$ & $V,\kappa\approx\mathrm{cst}$ \\
\bottomrule
\end{tabular}
\end{center}

\begin{derivation}[Padafhankelijkheid van arbeid]
Beschouw twee paden $A$ en $C$ van punt 1 naar punt 2. De cyclus $1\xrightarrow{C}2\xrightarrow{A}1$ geeft:
\[
  W_{\text{cyclus}} = W_{21}(C) + W_{12}(A)
\]
Omdat de oppervlaktes onder de twee curves in het $PV$-diagram verschillen, volgt
$|W_{21}(C)| \neq |W_{21}(A)|$ voor twee willekeurige paden, dus arbeid is padafhankelijk.
\end{derivation}

\subsection{Eerste Wet en Warmtecapaciteiten}

\begin{postulate}[Eerste Wet]
\[
  \Delta U = Q + W, \qquad dU = \dbar Q + \dbar W
\]
\end{postulate}

$U$ is een toestandsfunctie; $\dbar Q$ en $\dbar W$ zijn padafhankelijk.
Hydrostatisch: $dU = \dbar Q - P\,dV$. Overige systemen: zie tabel §1.3.

\[
  C_V = \pd[V]{\dbar Q}{T} = \pd[V]{U}{T} = T\pd[V]{S}{T},
  \qquad
  C_P = \pd[P]{\dbar Q}{T} = \pd[P]{U}{T} + P\pd[P]{V}{T} = T\pd[P]{S}{T}
\]

\subsection{Ideaal Gas}

Vrije expansie $\Rightarrow$ $(\partial U/\partial V)_T = 0$, dus $U = U(T)$ alleen en $C_V = dU/dT$.

\[
  \dbar Q = C_V\,dT + P\,dV = C_P\,dT - V\,dP
\]

\begin{derivation}[$C_P = C_V + nR$]
Differentieer $PV=nRT$: $P\,dV + V\,dP = nR\,dT$. Substitueer $P\,dV = nR\,dT - V\,dP$
in $\dbar Q = C_V\,dT + P\,dV$:
\[
  \dbar Q = (C_V + nR)\,dT - V\,dP
  \implies \pd[P]{\dbar Q}{T} = C_P = C_V + nR
\]
\end{derivation}

Monoatomisch: $\gamma=5/3$;\quad diatomisch: $\gamma\approx7/5$.

\begin{derivation}[Adiabatisch quasistatisch: $PV^\gamma = \mathrm{cst}$]
$\dbar Q=0$ toegepast op de twee vormen van $\dbar Q$:
\[
  0 = C_V\,dT + P\,dV \quad\text{en}\quad 0 = C_P\,dT - V\,dP
  \implies \frac{dP}{P} = -\gamma\,\frac{dV}{V}
\]
Integreren: $\ln P = -\gamma\ln V + \mathrm{cst}$, dus $\boxed{PV^\gamma = \mathrm{cst}}$.
Combineren met $PV = nRT$ geeft ook $TV^{\gamma-1}=\mathrm{cst}$ en $TP^{(1-\gamma)/\gamma}=\mathrm{cst}$.\\
Arbeid: $W = \dfrac{P_fV_f - P_iV_i}{\gamma-1} = C_V(T_f-T_i)$.
\end{derivation}

Methode van Rüchardt: $\gamma = 4\pi^2 mV/(\tau^2 P A^2)$.

%%=============================================================
\section{Tweede Wet en Thermische Machines (H6)}
%%=============================================================

\subsection{Rendement en Prestatiecoëfficiënt}

Eerste wet voor machines: $|Q_H| = |Q_K| + |W|$.
\[
  \eta = \frac{|W|}{|Q_H|} = 1 - \frac{|Q_K|}{|Q_H|},
  \qquad
  \omega_{\text{koel}} = \frac{|Q_K|}{|W|}
\]
\[
  \boxed{\eta_{\mathrm{Carnot}} = 1 - \frac{T_K}{T_H}}
  \qquad\text{(bovengrens voor alle machines)}
\]

\begin{center}
\renewcommand{\arraystretch}{1.2}
\begin{tabular}{@{}llll@{}}
\toprule
Cyclus   & Processen                           & Rendement             & Nota \\
\midrule
Carnot   & 2 isoth.\ + 2 adiab.                & $1 - T_K/T_H$         & reversibel \\
Stirling & 2 isoth.\ + 2 isoch.                & $= \eta_{\mathrm{Carnot}}$ & met regenerator \\
Otto     & 2 adiab.\ + 2 isoch.                & $1 - r^{1-\gamma}$    & $r=V_1/V_2$ \\
Diesel   & 2 adiab.\ + 1 isobaar + 1 isoch.    & zie afleiding         & $r_C=V_1/V_2$, $r_E=V_1/V_3$ \\
Rankine  & 6 proc.\ (stoom)                    & ---                   & zie cursus \\
\bottomrule
\end{tabular}
\end{center}

\begin{derivation}[Otto-rendement]
Warmte-uitwisseling enkel bij $2\to3$ (isochoor) en $4\to1$ (isochoor):
\[
  |Q_H| = C_V(T_3-T_2), \quad |Q_K| = C_V(T_4-T_1)
  \implies \eta = 1 - \frac{T_4-T_1}{T_3-T_2}
\]
Adiabatische relaties $1\to2$ en $3\to4$:
$T_1V_1^{\gamma-1}=T_2V_2^{\gamma-1}$ en $T_4V_1^{\gamma-1}=T_3V_2^{\gamma-1}$,
zodat $(T_4-T_1)/(T_3-T_2) = (V_2/V_1)^{\gamma-1} = r^{1-\gamma}$. Dus:
\[
  \boxed{\eta_{\mathrm{Otto}} = 1 - r^{1-\gamma}}
\]
\end{derivation}

\begin{derivation}[Diesel-rendement]
Warmte bij $2\to3$ (isobaar): $|Q_H|=C_P(T_3-T_2)$; bij $4\to1$ (isochoor): $|Q_K|=C_V(T_4-T_1)$.
\[
  \eta = 1 - \frac{T_4-T_1}{\gamma(T_3-T_2)}
\]
Adiabatische relaties en isobare relatie $T_3/T_2 = V_3/V_2$
elimineren de temperaturen ten voordele van $r_C, r_E$:
\[
  \boxed{\eta_{\mathrm{Diesel}} = 1 - \frac{1}{\gamma}\,\frac{r_E^{-\gamma}-r_C^{-\gamma}}{r_E^{-1}-r_C^{-1}}}
\]
\end{derivation}

\subsection{Tweede Wet: Formuleringen en Equivalentie}

\begin{postulate}[Kelvin-Planck]
Het enige resultaat van een cyclus kan niet zijn dat warmte van één warme bron
volledig in arbeid wordt omgezet.
\end{postulate}
\begin{postulate}[Clausius]
Het enige resultaat van een cyclus kan niet zijn dat warmte van koud naar warm stroomt.
\end{postulate}

\begin{derivation}[Equivalentie K-P $\Leftrightarrow$ Clausius]
\textbf{K $\Rightarrow$ C}: Stel een machine schendt Clausius: ze transporteert $Q_2$ van koud naar warm
zonder arbeid. Koppel er een Kelvin-Planck-machine aan die $W = Q_1$ omzet in arbeid en $Q_2$
afgeeft aan het koude reservoir. Netto: $Q_1 + Q_2$ onttrokken aan warm, $Q_2$ naar koud, netto arbeid $Q_1$. Dit schendt K-P. \\[2pt]
\textbf{C $\Rightarrow$ K}: Stel een machine schendt K-P: ze zet $Q_1$ volledig om in arbeid. Gebruik die arbeid voor een warmtepomp die $Q_2$ van koud naar warm pompt.
Netto: $Q_2$ stroomt van koud naar warm, wat Clausius schendt. \hfill$\square$
\end{derivation}

%%=============================================================
\section{Reversibiliteit, Entropie en Derde Wet (H7--H8)}
%%=============================================================

\subsection{Reversibiliteit en RA-oppervlakken}

\textbf{Reversibel} $\Leftrightarrow$ quasistatisch én geen dissipatieve effecten.

\begin{postulate}[Carathéodory]
In de buurt van elke evenwichtstoestand zijn er toestanden die via een adiabatisch
proces niet bereikbaar zijn.
\end{postulate}

Algemeen systeem met veralgemeende kracht $Y$ en verplaatsing $X$:
$\dbar Q = dU - Y\,dX$. Bij $\dbar Q=0$:
\[
  \frac{dT_e}{dX} = \frac{Y - (\partial U/\partial X)_{T_e}}{(\partial U/\partial T_e)_X}
  \implies \sigma(T_e, X) = \mathrm{cst}
  \quad\text{(RA-oppervlak)}
\]
RA-oppervlakken snijden elkaar nooit; dimensie = toestandsruimte $-1$.

\subsection{Integreerbaarheid van $\dbar Q$ en Kelvinschaal}

\begin{derivation}[Integrerende factor en entropie]
Schrijf $U = U(\sigma, X, X')$. Vergelijk $dU$ met $\dbar Q = dU - Y\,dX - Y'\,dX'$:
\[
  \dbar Q = \pd[\sigma,X']{U}{\sigma}\,d\sigma
  + \Bigl[\pd[\sigma,X']{U}{X}-Y\Bigr]dX
  + \Bigl[\pd[\sigma,X]{U}{X'}-Y'\Bigr]dX' = 0
\]
Doordat $X, X'$ onafhankelijk zijn: $(\partial U/\partial X)_{\sigma,X'}=Y$ en
$(\partial U/\partial X')_{\sigma,X}=Y'$, zodat $\dbar Q = \lambda(\sigma)\,d\sigma$
met $\lambda(\sigma) = (\partial U/\partial\sigma)_{X,X'}$.\\[4pt]
Via analyse van twee gekoppelde systemen volgt dat de integrerende factor de vorm
$1/\phi(T_e)$ heeft met $\phi$ uitsluitend afhankelijk van $T_e$:
\[
  \dbar Q = \phi(T_e)\,f(\sigma)\,d\sigma
\]
\end{derivation}

\begin{derivation}[Kelvinschaal = ideale gastemperatuur]
Carnot-cyclus op ideaal gas: isotherm $1\to2$ bij $T$ en $3\to4$ bij $T'$:
\[
  Q = nRT\ln\!\tfrac{V_2}{V_1}, \quad Q' = nRT'\ln\!\tfrac{V_4}{V_3}
\]
Adiabaten $2\to3$ en $4\to1$: $C_V\,dT = -(nRT/V)\,dV$ $\Rightarrow$
$\ln(V_2/V_3) = \ln(V_1/V_4)$ $\Rightarrow$ $V_2/V_1 = V_4/V_3$. Dus:
\[
  \frac{Q}{Q'} = \frac{T}{T'} = \frac{T_k}{T_k'}
\]
De Kelvintemperatuur $T_k$ is stof-onafhankelijk en gelijk aan de ideale gastemperatuur. \hfill$\square$
\end{derivation}

\begin{postulate}[Derde Wet (Nernst)]
Het absolute nulpunt $T=0$ is met een eindig aantal stappen onbereikbaar.
\end{postulate}

\subsection{Entropie: definitie en eigenschappen}

\[
  \boxed{dS = \frac{\dbar Q_{\mathrm{rev}}}{T}},
  \qquad
  \Delta S = \int_{i,\,\mathrm{rev}}^{f}\frac{\dbar Q}{T}
\]

Enkel $\Delta S$ is gedefinieerd.
Reversibele cyclus: $\oint \dbar Q_{\mathrm{rev}}/T = 0$ (Clausius-theorema).

\begin{derivation}[Entropie ideaal gas]
Uit $\dbar Q = C_V\,dT + P\,dV$ en $P/T = nR/V$:
\[
  \frac{\dbar Q}{T} = \frac{C_V}{T}\,dT + \frac{nR}{V}\,dV
  \implies \boxed{S = C_V\ln T + nR\ln V + S_0} \quad (C_V=\mathrm{cst})
\]
Analoog via $\dbar Q = C_P\,dT - V\,dP$ en ideale gaswet:
$S = C_P\ln T - nR\ln V + S_0$.
\end{derivation}

Isentroop $\equiv$ reversibel adiabatisch ($dS=0$).
TS-diagram: $Q_R = \int T\,dS$ (oppervlakte onder curve).

\subsection{Entropieprincipe}

\begin{derivation}[$\Delta S_{\mathrm{univ}} \geq 0$]
Beschouw een irreversibel adiabatisch proces $i\to f$ ($\Delta S = S_f - S_i$),
gevolgd door een irreversibel adiabatisch proc. $f\to k$ en een reversibel
isotherm proc. $k\to j$ bij $T'$, en tenslotte een reversibel adiabatisch proc.
$j\to i$. Netto entropie van de cyclus = 0:
\[
  \Delta S + (S_j - S_k) = 0 \implies \Delta S = S_k - S_j
\]
De enige warmtestroom is $Q_R = T'(S_j - S_k)$ bij $k\to j$.
Tweede wet: $Q_R \leq 0$, dus $T'(S_j-S_k)\leq 0$, dus $\Delta S = S_k - S_j \geq 0$.
Zou $\Delta S = 0$, dan zou het proces reversibel zijn: tegenspraak. Dus $\Delta S > 0$. \hfill$\square$
\end{derivation}

\begin{postulate}[Entropieprincipe]
$\Delta S_{\mathrm{univ}} \geq 0$, met gelijkheid enkel voor reversibele processen.
\end{postulate}

Bijzondere entropieveranderingen:
\begin{center}
\renewcommand{\arraystretch}{1.2}
\begin{tabular}{@{}ll@{}}
\toprule
Irreversibel proces & $\Delta S$ \\
\midrule
Uitw.\ mechanisch (isobaar vervanging)   & $C_P\ln(T_f/T_i)$ \\
Int.\ mech.\ (isotherm ideaal gas)       & $nR\ln(V_f/V_i)$ \\
Externe thermische                        & $\Delta S_{\mathrm{univ}} = Q/T_K - Q/T_H$ \\
\bottomrule
\end{tabular}
\end{center}

\subsection{Machineprestaties en Niet-beschikbare Energie}

\begin{derivation}[$W_{\max}$ en $E = T_0\,\Delta S$]
$\Delta S_{\mathrm{univ}} = (Q_H-W)/T_K - Q_H/T_H \geq 0$
$\Rightarrow$ $W \leq Q_H(1 - T_K/T_H)$, zodat:
\[
  W_{\max} = Q_H\!\left(1 - \frac{T_K}{T_H}\right)
\]
Na een irreversibel proces $i\to j$ in contact met reservoir $T_0$:
reversibel herstelproces voert warmte $Q_{\mathrm{rev}} = T_0(S_f-S_i)$ af.
Die energie is niet in arbeid om te zetten:
\[
  \boxed{E = T_0\,\Delta S} \quad\text{(niet-beschikbare energie)}
\]
\end{derivation}

\[
  W_{\mathrm{min,koel}} = (S_1-S_2)T_K - Q
\]

%%=============================================================
\section{Thermodynamische Potentialen en Relaties (H9)}
%%=============================================================

\subsection{De Vier Thermodynamische Potentialen}

\begin{center}
\renewcommand{\arraystretch}{1.3}
\begin{tabular}{@{}lllll@{}}
\toprule
Potentiaal       & Definitie & Differentiaalvorm     & Nat.\ var. & Gebruik \\
\midrule
$U$ (intern)     & ---       & $dU = T\,dS - P\,dV$  & $S,V$      & algemeen \\
$H = U+PV$       & enthalpie & $dH = T\,dS + V\,dP$  & $S,P$      & isobaar; $Q=\Delta H$ \\
$F = U-TS$       & Helmholtz & $dF = -S\,dT - P\,dV$ & $T,V$      & isoth.--isoch.; $F=\mathrm{cst}$ \\
$G = H-TS$       & Gibbs     & $dG = -S\,dT + V\,dP$ & $T,P$      & isoth.--isobaar; $G=\mathrm{cst}$ \\
\bottomrule
\end{tabular}
\end{center}

Partiële afgeleiden:
\[
  \pd[P]{H}{S}=T,\quad\pd[S]{H}{P}=V;\qquad
  \pd[V]{F}{T}=-S,\quad\pd[T]{F}{V}=-P;\qquad
  \pd[P]{G}{T}=-S,\quad\pd[T]{G}{P}=V
\]

\textbf{Enthalpie}: bij isobaar reversibel proc.: $H_f - H_i = \int C_P\,dT = Q$.\\
\textbf{Smoorproces} (Joule-Kelvin, irreversibel, $H_i = H_f$):
\[
  \mu = \pd[H]{T}{P} = \frac{1}{C_P}\!\left[T\pd[P]{V}{T} - V\right] = \frac{V(T\beta-1)}{C_P}
  \qquad (\mu=0 \text{ voor ideaal gas})
\]

\textbf{Gibbsfunctie bij fase-evenwichten}:
$G_{\text{vast}} = G_{\text{vl.}}$ (smelt),\;
$G_{\text{vl.}} = G_{\text{damp}}$ (verdamp.),\;
$G_{\text{vast}} = G_{\text{damp}}$ (sublim.)

\subsection{Maxwell-betrekkingen}

Totale differentiaal $df = M\,dx + N\,dy$ $\Rightarrow$ $(\partial M/\partial y)_x = (\partial N/\partial x)_y$.

\begin{derivation}[Maxwell-betrekkingen]
Pas de symmetrie van gemengde tweede afgeleiden toe op elk potentiaal:
\begin{center}
\renewcommand{\arraystretch}{1.4}
\begin{tabular}{@{}ll@{}}
\toprule
Uit $dU = T\,dS - P\,dV$  & $\pd[S]{T}{V} = -\pd[V]{P}{S}$ \\
Uit $dH = T\,dS + V\,dP$  & $\pd[S]{T}{P} = \pd[P]{V}{S}$ \\
Uit $dF = -S\,dT - P\,dV$ & $\pd[T]{S}{V} = \pd[V]{P}{T}$ \\
Uit $dG = -S\,dT + V\,dP$ & $\pd[T]{S}{P} = -\pd[P]{V}{T}$ \\
\bottomrule
\end{tabular}
\end{center}
\end{derivation}

\subsection{TdS-vergelijkingen}

\begin{derivation}[TdS-vergelijkingen]
Schrijf $S = S(T,V)$:\;
$dS = (\partial S/\partial T)_V\,dT + (\partial S/\partial V)_T\,dV$.
Isochoor: $C_V = T(\partial S/\partial T)_V$.
Maxwell $dF$: $(\partial S/\partial V)_T = (\partial P/\partial T)_V$:
\[
  \boxed{TdS = C_V\,dT + T\pd[V]{P}{T}\,dV} \tag{TdS-1}
\]
Schrijf $S = S(T,P)$:\;
isotherm: $C_P = T(\partial S/\partial T)_P$.
Maxwell $dG$: $(\partial S/\partial P)_T = -(\partial V/\partial T)_P$:
\[
  \boxed{TdS = C_P\,dT - T\pd[P]{V}{T}\,dP} \tag{TdS-2}
\]
Schrijf $S = S(P,V)$ en gebruik kettingregel + $C_V$ en $C_P$:
\[
  \boxed{TdS = C_V\pd[V]{T}{P}\,dP + C_P\pd[P]{T}{V}\,dV} \tag{TdS-3}
\]
\end{derivation}

Toepassingen TdS-vergelijkingen:
\begin{itemize}
  \item \textbf{Rev.\ isotherm}: $Q = -T\!\int V\beta\,dP \approx -T\,\overline{V\beta}(P_f-P_i)$
  \item \textbf{Rev.\ adiabatisch}: $dT = (TV\beta/C_P)\,dP$
\end{itemize}

\subsection{Energie- en Enthalpievergelijkingen}

\begin{derivation}[Energie- en enthalpievergelijkingen]
Uit $dU = T\,dS - P\,dV$, deel door $dV$ (bij $T=\mathrm{cst}$) en gebruik
Maxwell $(\partial S/\partial V)_T = (\partial P/\partial T)_V$:
\[
  \pd[T]{U}{V} = T\pd[V]{P}{T} - P \quad\text{(1e energievgl.)}
\]
Deel $dU$ door $dP$ (bij $T=\mathrm{cst}$) en gebruik Maxwell $(\partial S/\partial P)_T = -(\partial V/\partial T)_P$:
\[
  \pd[T]{U}{P} = -P\pd[T]{V}{P} \quad\text{(2e energievgl.)}
\]
Analoog voor $dH = T\,dS + V\,dP$:
\[
  \pd[T]{H}{P} = T\pd[T]{S}{P} + V \quad\text{(1e enthalpievgl.)}
\]
\[
  \pd[T]{H}{V} = T\pd[T]{S}{V} + V\pd[T]{P}{V} \quad\text{(2e enthalpievgl.)}
\]
\end{derivation}

\subsection{Warmtecapaciteitsvergelijkingen}

\begin{derivation}[$C_P - C_V$]
Stel TdS-1 = TdS-2 en los op naar $dT$ (met $T=T(P,V)$):
\[
  dT = \frac{T\pd[P]{V}{T}}{C_P-C_V}\,dP + \frac{T\pd[V]{P}{T}}{C_P-C_V}\,dV
  = \pd[V]{T}{P}\,dP + \pd[P]{T}{V}\,dV
\]
Vergelijken geeft $C_P - C_V = T(\partial P/\partial T)_V(\partial V/\partial T)_P$.
Met de Jacobiaan-identiteit $(\partial P/\partial T)_V(\partial T/\partial V)_P(\partial V/\partial P)_T = -1$:
\[
  \boxed{C_P - C_V = -T\!\left(\pd[V]{P}{T}\right)^{\!2}\!\pd[T]{V}{P} = \frac{TV\beta^2}{\kappa}}
\]
\end{derivation}

Gevolgen: $C_P > C_V$ altijd; $C_P \to C_V$ als $T\to 0$ of $(\partial V/\partial T)_P = 0$.

\begin{derivation}[$C_P/C_V = \kappa/\kappa_S$]
Gebruik TdS-1 en TdS-2 bij $dS=0$:
\[
  C_V\,dT = -T\pd[V]{P}{T}\,dV_S, \quad C_P\,dT = T\pd[P]{V}{T}\,dP_S
\]
\[
  \frac{C_P}{C_V}
  = \frac{T(\partial V/\partial T)_P}{-T(\partial P/\partial T)_V}\,\pd[S]{P}{V}
  = -\frac{(\partial P/\partial V)_S}{(\partial P/\partial V)_T}
  = \boxed{\frac{\kappa}{\kappa_S}}
  \quad\text{met}\quad \kappa_S = -\frac{1}{V}\pd[S]{V}{P}
\]
\end{derivation}

\subsection{Fase-overgangen}

\subsubsection*{Eerste-orde (Clapeyron)}

$T$ en $P$ constant; $S$ en $V$ springen. Latente warmte: $L = T(s_f-s_i)$.
$C_P,\,\beta,\,\kappa \to \infty$ tijdens de overgang.

\begin{derivation}[Vergelijking van Clapeyron]
Langs de overgangskromme: $G_i = G_f$, dus bij een kleine stap $dG_i = dG_f$:
\[
  V_i\,dP - S_i\,dT = V_f\,dP - S_f\,dT
  \implies
  \boxed{\frac{dP}{dT} = \frac{S_f-S_i}{V_f-V_i} = \frac{L}{T(V_f-V_i)}}
\]
Alternatief via TdS-1 met $(\partial P/\partial T)_V = dP/dT$ (want $P=P(T)$ langs de kromme):
$TdS = T(dP/dT)(V_f-V_i) = L$. \hfill$\square$
\end{derivation}

Teken: positief als $V_f > V_i$ (expansie); negatief voor water bij smelting.

\subsubsection*{Tweede-orde (Ehrenfest)}

$G$, $S$, $V$ continu; $C_P$, $\beta$, $\kappa$ discontinu.
Voorbeeld: supergeleider $\to$ normaal geleidend.

\begin{derivation}[Vergelijkingen van Ehrenfest]
\textbf{1e vgl.:} Behoud van entropie ($S_i = S_f$) langs de overgangskromme,
dus $dS_i = dS_f$. Gebruik TdS-2 (bij $dS=0$ en isotherme variatie $dT,dP$ langs de kromme):
\[
  C_{P,i}\,dT - TV\beta_i\,dP = C_{P,f}\,dT - TV\beta_f\,dP
  \implies \boxed{\frac{dP}{dT} = \frac{\Delta C_P}{TV\,\Delta\beta}}
\]
\textbf{2e vgl.:} Behoud van volume ($V_i = V_f$), dus $dV_i = dV_f$:
\[
  \pd[P]{V}{T}\,dT + \pd[T]{V}{P}\,dP = \text{idem voor }f
  \implies \beta_i\,dT - \kappa_i\,dP = \beta_f\,dT - \kappa_f\,dP
  \implies \boxed{\frac{dP}{dT} = \frac{\Delta\beta}{\Delta\kappa}}
\]
\end{derivation}

\textbf{Lambda-transitie}: $C_P\to\infty$ maar $\beta,\kappa$ eindig $\Rightarrow$
Ehrenfest niet geldig ($T,P,G,S,V$ constant; $C_P,\beta,\kappa\to\infty$).

\end{document}